\chapter{What Is Reinforcement Learning?}
\label{ch:what-is-rl}

\section{Motivation: Why Does RL Exist?}

Suppose you want to teach a robot arm to pour a cup of coffee without
spilling it. Or teach a program to play chess well enough to beat a
grandmaster. Or teach a recommendation engine to keep a user engaged over
weeks, not just the next click. What do these three problems have in
common, and why does neither ordinary supervised learning nor unsupervised
learning solve them well?

Think about what supervised learning needs: a labeled dataset of
\emph{correct answers}. For coffee-pouring, that would mean someone
already knows, for every possible arm configuration, exactly which torque
to apply next — but that's the entire problem we're trying to solve in
the first place. There is no ``correct label'' lying around for
\emph{sequential decision making}, because the right decision today
depends on decisions you haven't made yet, and on consequences that
unfold over time.

\begin{intuitionbox}
Supervised learning answers the question \emph{``what is this?''}
Reinforcement learning answers the much harder question
\emph{``what should I do, right now, given that my choice affects
everything that happens afterward?''}
\end{intuitionbox}

This is precisely the gap reinforcement learning (RL) was built to fill:
a mathematical framework, and a family of algorithms, for learning
\emph{good sequences of decisions} through trial, feedback, and
experience — without anyone needing to hand-label the correct action at
every step.

\section{RL vs.\ Supervised Learning vs.\ Unsupervised Learning}

It's worth being precise about the differences, because the vocabulary of
machine learning ("training," "loss," "data") gets reused across all
three paradigms in ways that can blur what is actually different.

\begin{center}
\begin{tabular}{@{}p{2.6cm}p{3.6cm}p{3.6cm}p{3.6cm}@{}}
\toprule
& \textbf{Supervised Learning} & \textbf{Unsupervised Learning} & \textbf{Reinforcement Learning} \\
\midrule
\textbf{Data} & Labeled examples $(x, y)$ & Unlabeled examples $x$ & Experience: sequences of $(\state, \act, \rew)$ \\
\textbf{Feedback} & Correct answer given directly & No explicit feedback & A scalar reward, often delayed \\
\textbf{Goal} & Predict $y$ from $x$ & Discover structure in $x$ & Choose actions that maximize \emph{cumulative} future reward \\
\textbf{Data source} & Fixed dataset, i.i.d.\ typically & Fixed dataset & Data is generated by the agent's own actions \\
\textbf{Time} & Usually no temporal structure & Usually no temporal structure & Fundamentally sequential \\
\bottomrule
\end{tabular}
\end{center}

\begin{keyidea}
The single most important difference is this: in supervised learning, the
data is handed to you. In reinforcement learning, the agent's actions
\emph{determine what data it sees next}. Choosing badly doesn't just cost
you a wrong label — it can trap you in a part of the world where you never
observe the good outcomes at all. This is why exploration becomes a
central, unavoidable concern in RL (we'll return to this properly in
Chapter~26), and it has no real analogue in supervised learning.
\end{keyidea}

\begin{commonmistakes}
A common early misconception is to think of RL as ``supervised learning,
but the labels are rewards instead of classes.'' This undersells the
difference. A reward is not a label telling you what the right action
was — it's a scalar judgment of how good an outcome turned out to be,
often received many steps after the action that caused it, and the agent
must itself figure out \emph{which} of its past actions deserves credit
for it. That credit-assignment problem doesn't exist in supervised
learning at all.
\end{commonmistakes}

\section{RL as Sequential Decision Making}

The right way to picture an RL problem is as a loop, repeating over
discrete time steps $t = 0, 1, 2, \dots$:

\begin{enumerate}
  \item The agent observes the current state of the world.
  \item Based on that state, the agent chooses an action.
  \item The world changes in response to that action, transitioning to a
        new state.
  \item The agent receives a reward — a single number — indicating how
        good or bad that transition was.
  \item The loop repeats, with the agent now facing a new state.
\end{enumerate}

This loop is the single most important picture in this entire book.
Nearly everything from here through Chapter~61 is either (a) a more
precise mathematical description of some piece of this loop, or (b) an
algorithm for making the agent better at choosing actions inside it.

\begin{center}
\begin{tikzpicture}[
    node distance=2.4cm,
    box/.style={draw, thick, rounded corners, minimum width=3.2cm,
                minimum height=1.4cm, align=center, fill=blue!6},
    every node/.style={font=\small}
]
\node[box] (agent) {Agent};
\node[box, right=of agent, fill=orange!8] (env) {Environment};

\draw[-{Latex[length=2.5mm]}] ($(agent.north east)+(0,0.35)$) -- ($(env.north west)+(0,0.35)$)
    node[midway, above] {action $\act_t$};
\draw[-{Latex[length=2.5mm]}] ($(env.south west)+(0,-0.35)$) -- ($(agent.south east)+(0,-0.35)$)
    node[midway, below] {state $\state_{t+1}$, reward $\rew_{t+1}$};
\end{tikzpicture}
\end{center}

\begin{intuitionbox}
Picture a thermostat-controlled learner instead of a robot, if that's more
concrete: the \emph{agent} is the decision-maker (the thermostat's
controller); the \emph{environment} is everything outside the
agent's direct control (the room, the weather, the people opening
windows). The agent can only affect the world through its actions, and
can only learn about the world through the states and rewards it
receives back.
\end{intuitionbox}

\section{The Vocabulary of RL}

We now introduce the core terms precisely. Each one will reappear, in
increasingly formal clothing, throughout the rest of the book — this is
the first, informal pass.

\subsection*{Agent}
The learner and decision-maker. In our examples: the navigator choosing
moves in Grid World, the controller choosing torques for the Robotic Arm.

\subsection*{Environment}
Everything the agent interacts with but does not directly control.
The environment defines what happens (what new state results, what
reward is given) in response to the agent's action.

\subsection*{State}
\begin{definition}[State, informal]
The state $\state_t \in \stateSet$ is a description of the situation the
agent finds itself in at time $t$ — in principle, everything relevant to
predicting what happens next.
\end{definition}
In Grid World, the state is simply the agent's $(row, column)$ position.
In the Robotic Arm, the true state might include every joint angle and
joint velocity.

\subsection*{Observation}
\begin{definition}[Observation, informal]
The observation $\obs_t \in \obsSet$ is what the agent actually
\emph{perceives} at time $t$ — which may be the full state, or only a
partial, noisy, or processed view of it.
\end{definition}

\begin{whythisworks}
Why bother distinguishing state from observation at all? Because in many
real problems — including almost all of robotics — the agent never has
direct access to the true state. A camera image is an observation, not
the state; it doesn't directly tell you joint velocities, friction
coefficients, or what's occluded behind the gripper. For much of this
book (Chapters 1–37) we will simplify by assuming the agent's observation
\emph{is} the full state — this special case is important enough to have
its own name (a Markov Decision Process, coming in Chapter~4). We will
deliberately drop this simplifying assumption later, in Chapter~38
(POMDPs), once you have enough machinery to handle the general case.
\end{whythisworks}

\begin{center}
\begin{tabular}{@{}p{4.2cm}p{4.2cm}p{4.6cm}@{}}
\toprule
& \textbf{State $\state_t$} & \textbf{Observation $\obs_t$} \\
\midrule
What it is & The true, complete situation & What the agent actually gets to see \\
Grid World example & Exact $(row, col)$ cell & (Same, in this simple case) \\
Robotic Arm example & All joint angles/velocities, exact object pose & A camera image, or noisy joint sensors \\
Chapters that assume $\obs_t = \state_t$ & — & Ch.~1–37 (the ``fully observable'' assumption) \\
\bottomrule
\end{tabular}
\end{center}

\subsection*{Action}
\begin{definition}[Action, informal]
The action $\act_t \in \actSet$ is the choice the agent makes at time $t$,
selected from the set of actions available to it.
\end{definition}
Actions can be \emph{discrete} (Grid World: \{UP, DOWN, LEFT, RIGHT\}) or
\emph{continuous} (Robotic Arm: a real-valued torque per joint). We will
treat discrete actions almost exclusively until Chapter~22, because the
mathematics is simpler and the ideas transfer.

\subsection*{Reward}
\begin{definition}[Reward, informal]
The reward $\rew_{t+1}$ is a single scalar number, given to the agent by
the environment immediately after it takes action $\act_t$ in state
$\state_t$, indicating how good or bad that particular transition was.
\end{definition}

\begin{commonmistakes}
Beginners very often conflate \emph{reward} with \emph{value} or with
\emph{return}. A reward is a single, immediate number for a single step.
It says nothing, by itself, about long-term success. An action can have a
great immediate reward and still be a terrible choice overall (walking
off the cliff in Cliff Walking might, in principle, offer a shortcut with
small per-step costs before disaster — the point is that reward is local,
not cumulative). We will formally define \emph{return} (the discounted
sum of future rewards) in Chapter~5, and \emph{value} (the expected
return) in Chapter~6. Keep these three words — reward, return, value —
mentally separate starting now; conflating them is the single most common
beginner error in this entire subject.
\end{commonmistakes}

\subsection*{Episode and Trajectory}
\begin{definition}[Episode]
An episode is one complete run of agent-environment interaction, from an
initial state to a terminal state (e.g., reaching the goal, falling off
the cliff, or hitting a maximum time limit).
\end{definition}

\begin{definition}[Trajectory]
A trajectory $\traj$ is the recorded sequence of states, actions, and
rewards experienced during (all or part of) an episode:
$$
\traj = (\state_0, \act_0, \rew_1, \state_1, \act_1, \rew_2, \state_2, \dots)
$$
\end{definition}

Some tasks are naturally \emph{episodic} (Grid World: the episode ends
when the goal is reached), while others are \emph{continuing}, with no
natural endpoint (e.g., a thermostat that runs forever). We will handle
both carefully once we define return in Chapter~5 — the discount factor
$\discount$ turns out to be exactly the tool that makes continuing tasks
mathematically well-behaved.

\section{A Complete, Worked Example: One Episode in Grid World}

Let's make all of this concrete with actual numbers, using Example A
from Chapter~0.

\begin{example}[One Grid World episode]
Consider a $1 \times 4$ strip of cells, labeled $0, 1, 2, 3$, where cell
$3$ is the goal. The agent starts in cell $0$. Actions are LEFT or RIGHT;
moving off the strip leaves the agent in place. Every step costs
$\rew = -1$, except the step that reaches the goal, which gives
$\rew = +10$.

One possible trajectory:
\begin{center}
\begin{tabular}{@{}cccc@{}}
\toprule
Step $t$ & State $\state_t$ & Action $\act_t$ & Reward $\rew_{t+1}$ \\
\midrule
0 & 0 & RIGHT & $-1$ \\
1 & 1 & RIGHT & $-1$ \\
2 & 2 & RIGHT & $+10$ \\
\bottomrule
\end{tabular}
\end{center}
so $\state_3 = 3$ (the goal) and the episode ends. Written as a
trajectory:
$$
\traj = (0, \text{RIGHT}, -1,\; 1, \text{RIGHT}, -1,\; 2, \text{RIGHT}, +10,\; 3).
$$
\end{example}

This trajectory already lets us ask the central question of the whole
book, informally: \emph{was RIGHT, RIGHT, RIGHT a good sequence of
actions?} Clearly yes here — it reached the goal in the minimum possible
number of steps. But notice we can't judge \emph{any single action} in
isolation from the others; the first RIGHT only ``paid off'' because of
what came after it. This is exactly the sequential credit-assignment
problem RL exists to solve, and it's why we'll need the concept of
\emph{return} (Chapter~5) before we can precisely define what ``good''
even means.

\section{Real-World Motivation}

RL is not an abstract curiosity — it is the right framework whenever a
decision-maker acts repeatedly, over time, and only finds out the
consequences of its choices later. A few concrete settings:

\begin{description}[leftmargin=1.6cm, style=nextline]
  \item[Robotics] \hfill \\
  A robot arm doesn't get a labeled "correct torque" — it gets to try, and
  observe whether the task succeeded. This is the central motivation for
  Part~X of this book.

  \item[Games] \hfill \\
  From Backgammon (one of RL's earliest successes) to Go and StarCraft II,
  games provide a clean, simulatable environment where an agent can play
  millions of episodes against itself.

  \item[Recommendation systems] \hfill \\
  What to show a user next affects not just this click, but their
  engagement over the following days — a genuinely sequential problem
  dressed up as a single-step prediction if you're not careful.

  \item[Autonomous systems] \hfill \\
  Self-driving cars, warehouse robots, and drones all make long sequences
  of interdependent decisions under uncertainty.
\end{description}

\begin{robotics}
Every vocabulary term introduced in this chapter will map directly onto
robotics as the book progresses: the \emph{state} becomes joint angles
and velocities (and eventually camera images); the \emph{action} becomes
a torque or target joint position; the \emph{reward} becomes however you
choose to score task success (Chapter~27 is entirely about how to do this
well); and the \emph{agent} becomes the control policy you are trying to
learn. Keep this mapping in mind — we will make it explicit again at the
start of Part~X.
\end{robotics}

\section{Limitations of This First Picture}

The agent-environment loop above is deliberately informal. It leaves
several questions unanswered, which is precisely why the next three
chapters exist:

\begin{itemize}
  \item \textbf{What mathematical assumptions make this loop tractable?}
  We've been assuming, implicitly, that the future only depends on the
  present state and action — not on the entire history. This assumption
  has a name (the \emph{Markov property}) and it is not automatic; we
  examine it carefully in Chapter~3.
  \item \textbf{How do we formally model transitions and rewards?}
  Chapter~4 introduces the Markov Decision Process (MDP), the
  five-part mathematical object $(\stateSet, \actSet, P, R, \discount)$
  that makes everything in this loop precise.
  \item \textbf{How do we judge a sequence of rewards as ``good'' or
  ``bad''?} We've informally judged our Grid World trajectory as good, but
  we have no formula yet. Chapter~5 introduces \emph{return}.
\end{itemize}

\begin{keyidea}
This chapter gave you the \emph{shape} of every RL problem: an agent and
an environment, interacting in a loop, generating states, actions, and
rewards over time. Everything else in this book is about making that
shape mathematically precise (Chapters 2–7), and then building
algorithms that let an agent get good at it, from small tabular
methods (Chapters 8–12) all the way to deep RL, continuous control, and
multi-agent systems.
\end{keyidea}

\section*{Exercises}
% Exercises for Chapter 1 — What Is Reinforcement Learning?

\begin{enumerate}
  \item \textbf{Conceptual.} Explain, in your own words, why supervised
  learning is not a natural fit for teaching a robot to pour coffee.
  What specifically is missing from a supervised-learning dataset that RL
  provides instead?

  \item \textbf{Conceptual.} Give an example (not from this chapter) of a
  real-world problem that is naturally episodic, and one that is
  naturally continuing. Justify each choice.

  \item \textbf{State vs.\ observation.} For a self-driving car, describe
  something that plausibly belongs to the true \emph{state} of the world
  but would \emph{not} be directly available in the car's
  \emph{observation} from its sensors.

  \item \textbf{Numerical.} Using the $1\times4$ Grid World from this
  chapter's worked example, write out the full trajectory (states,
  actions, rewards) if the agent instead takes the actions
  LEFT, RIGHT, RIGHT, RIGHT, RIGHT (recall moving off the strip leaves
  the agent in place). How many total steps does this trajectory take
  to reach the goal, compared to the RIGHT, RIGHT, RIGHT trajectory in
  the chapter?

  \item \textbf{Thought experiment.} Suppose every reward in Grid World
  were changed to exactly $0$, including at the goal. What, if anything,
  could an RL agent learn in this environment? What does this tell you
  about the role reward plays in defining the task itself?

  \item \textbf{Robotics.} Pick one component of the agent-environment
  loop (state, action, reward, or environment) and describe concretely
  what it would correspond to for a warehouse robot picking items off
  shelves. Be specific about units or data type (e.g., "action = a
  3-dimensional continuous vector of end-effector velocity, in m/s").

  \item \textbf{Common mistakes.} A classmate says: "Reward and value are
  basically the same thing, since a bigger reward is always better."
  Explain, using an example, why this statement is misleading.
\end{enumerate}

